% ============================================================================
% KAPITEL 14: Stationäre Kosmologie
% ============================================================================

\chapter{Stationäre Kosmologie}
\label{ch:kosmologie}

\section{Einleitung}
\label{sec:kosmo_einleitung}

Die WDBT+ führt zu einem stationären, nicht-expandierenden Universum. Die beobachtete Rotverschiebung ist keine Folge der Expansion des Raumes, sondern eine kumulative gravitative Wirkung auf Photonen während ihrer Reise durch das \( \Omega \)-Feld. Die DSTT-Drift (\( \dot{\beta} > 0 \)) bewirkt eine langsame Ausdehnung von Strukturen, die durch kontinuierliche Materieentstehung aus dem Quantenpotential kompensiert wird.

Die in diesem Kapitel entwickelte Skalenabhängigkeit der effektiven Hubble-Konstanten ist eng mit der Neutrinomasse verbunden: Beide Phänomene werden durch dieselbe Korrelationslänge \( L_{Q,0} \approx 2{,}0 \times 10^{46} \, \text{m} \) bestimmt, die aus der fraktalen Geometrie des Raumes folgt (siehe Anhang~\ref{app:neutrino_oszillation}). Dies ist ein zentrales konsistentes Element der Theorie.

\section{Grundlagen der Rotverschiebung in der WDBT}
\label{sec:kosmo_rotverschiebung}

\subsection{Allgemeine Form}
\label{sub:kosmo_rot_allgemein}

Die allgemeine, aus den Weber-Kräften folgende Rotverschiebung für ein Lichtsignal, das eine radialsymmetrische Massenverteilung \( M(r) \) durchläuft, lautet:

\[
z = \frac{1}{c^2} \int_0^d \frac{G M(r)}{r^2} \, dr \qquad \text{(allgemeine Form)}
\]

Diese Gleichung gilt für jede Massenverteilung und enthält keine fraktale Dimension \( D \). Sie ist die Grundlage der Rotverschiebung in der WDBT.

\subsection{Rotverschiebung ohne Expansion}
\label{sub:kosmo_rot_ohne_expansion}

In der WDBT+ ist die Rotverschiebung keine Doppler-Verschiebung durch Expansion, sondern eine kumulative gravitative Wirkung auf Photonen während ihrer Reise durch das \( \Omega \)-Feld (siehe Kapitel~\ref{ch:maxwell_gravitation}). Die allgemeine Form (siehe oben) beschreibt genau diesen Effekt.

\subsection{Fraktale Dichteverteilung als Spezialfall}
\label{sub:kosmo_rot_fraktal}

Nimmt man speziell die fraktale Dichteverteilung der WDBT+ an (siehe Kapitel~\ref{ch:fraktale_dimension}),

\[
\rho(r) = \rho_0 \left( \frac{r}{l_0} \right)^{D-3},
\]

so ergibt sich für die eingeschlossene Masse:

\[
M(r) = \int_0^r \rho(r') \cdot 4\pi r'^2 \, dr' = 4\pi \rho_0 l_0^{3-D} \frac{r^D}{D}.
\]

Einsetzen in die allgemeine Form liefert:

\[
z(d) = \frac{1}{c^2} \int_0^d \frac{G}{r^2} \cdot 4\pi \rho_0 l_0^{3-D} \frac{r^D}{D} \, dr
     = \frac{4\pi G \rho_0 l_0^{3-D}}{D c^2} \int_0^d r^{D-2} \, dr.
\]

Das Integral ergibt \( \int_0^d r^{D-2} dr = d^{D-1} / (D-1) \), und man erhält:

\[
z(d) = \frac{4\pi G \rho_0 l_0^{3-D}}{D(D-1) c^2} \, d^{\,D-1}.
\]

Dies ist ein Spezialfall der allgemeinen WDBT-Rotverschiebung für den Fall einer fraktalen Dichteverteilung. Für \( D \to 3 \) geht die fraktale Verteilung in eine konstante Dichte über.

\subsection{Die effektive Hubble-Konstante}
\label{sub:kosmo_effektive_hubble}

In der WDBT+ gibt es keine universelle Hubble-Konstante im Sinne der expandierenden Kosmologie, da das Universum stationär ist. Stattdessen definiert man eine \textbf{skalenabhängige effektive Hubble-Konstante}:

\[
H_{\text{eff}}(d) = c \cdot \frac{z(d)}{d}
\]

Aus der allgemeinen Rotverschiebungsformel (Abschnitt~\ref{sub:kosmo_rot_allgemein}) folgt für beliebige Massenverteilungen \( M(r) \):

\[
H_{\text{eff}}(d) = \frac{1}{d} \int_0^d \frac{GM(r)}{r^2} \, dr
\]

Für den Spezialfall der fraktalen Dichteverteilung (Abschnitt~\ref{sub:kosmo_rot_fraktal}) ergibt sich:

\[
H_{\text{eff}}(d) = \frac{4\pi G \rho_0 l_0^{3-D}}{D(D-1)c} \cdot d^{\,D-2}
\]

Mit \( D \approx 2,704 \) ist \( D-2 = 0,704 \), also:

\[
H_{\text{eff}}(d) \propto d^{0,704}
\]

Die Korrelationslänge des Q-Feldes ist \( L_{Q,0} \approx 2{,}0 \times 10^{46} \, \text{m} \). Sie ist nicht willkürlich, sondern folgt aus der fraktalen Geometrie des Raumes und der gemessenen Neutrinomasse (siehe Kapitel~\ref{ch:neutrino} und Anhang~\ref{app:neutrino_oszillation}). Dieselbe Skala \( L_{Q,0} \) bestimmt auch die Neutrinomasse:

\[
m_\nu = \frac{\hbar}{l_0 c} \left( \frac{l_0}{L_{Q,0}} \right)^{(3-D)/2} \approx 0,1 \, \text{eV}/c^2
\]

Damit sind die Neutrinomasse und die Skalenabhängigkeit der Hubble-Konstanten zwei Seiten derselben physikalischen Ursache – der fraktalen Struktur des Raumes.

\subsection{Die Hubble-Spannung als Skalenabhängigkeit}
\label{sub:kosmo_hubble_spannung}

Die Hubble-Spannung – die Diskrepanz zwischen verschiedenen Messmethoden für \( H_0 \) im \(\Lambda\)CDM-Modell – ist in der WDBT+ kein Rätsel, sondern eine Konsequenz der fraktalen Skalenabhängigkeit. Unterschiedliche Messmethoden operieren auf unterschiedlichen effektiven Skalen \( d \) und müssen daher systematisch unterschiedliche Werte für \( H_{\text{eff}} \) liefern.

Bezeichnet man mit \( H_0 \) den Wert der effektiven Hubble-Konstanten bei einer Referenzskala \( d_0 \) (z. B. bei lokalen Entfernungsmessungen), so gilt:

\[
H_{\text{eff}}(d) = H_0 \left( \frac{d}{d_0} \right)^{0,704}
\]

Die Hubble-Spannung lässt sich nun wie folgt erklären:

\[
\frac{H_{\text{local}}}{H_{\text{CMB}}} = \left( \frac{d_{\text{local}}}{d_{\text{CMB}}} \right)^{0,704}
\]

Mit den gemessenen Werten \( H_{\text{CMB}} \approx 67,4 \, \text{km/s/Mpc} \) und \( H_{\text{local}} \approx 73,5 \, \text{km/s/Mpc} \) ergibt sich:

\[
\frac{d_{\text{local}}}{d_{\text{CMB}}} = \left( \frac{73,5}{67,4} \right)^{1/0,704} \approx 1,129
\]

Die lokalen Messungen operieren also auf einer um etwa 12,9 \% \textbf{größeren} effektiven Skala als die CMB-Messungen (\( d_{\text{local}} > d_{\text{CMB}} \)).

Da der Exponent \( D-2 = 0,704 \) positiv ist, ist \( H_{\text{eff}}(d) \) eine monoton wachsende Funktion von \( d \). Aus \( d_{\text{local}} > d_{\text{CMB}} \) folgt daher:

\[
H_{\text{eff}}(d_{\text{local}}) > H_{\text{eff}}(d_{\text{CMB}})
\]

also:

\[
H_{\text{local}} > H_{\text{CMB}}
\]

Dies entspricht exakt der beobachteten Hubble-Spannung (\( 73,5 > 67,4 \)).

In einem stationären Universum ist dies kein Zeichen für eine Expansion oder eine Urknall-Epoche, sondern eine direkte Konsequenz der fraktalen Geometrie des Raumes.

Da \( \beta \) für Photonen von der Energie abhängt, zeigen verschiedene Frequenzen leicht unterschiedliche Rotverschiebungen. Dieser Effekt ist klein, aber mit Pulsar-Laufzeiten oder hochpräzisen Spektren messbar.

\section{Galaxienrotation ohne Dunkle Materie}
\label{sec:kosmo_galaxienrotation}

Flache Rotationskurven werden in der WDBT+ durch das \( \Omega \)-Feld und die anisotrope Trägheit der DSTT erklärt (siehe Kapitel~\ref{ch:dstt_weber_kraefte} und \ref{ch:maxwell_gravitation}). Die radiale Dichteverteilung einer Galaxie ergibt sich aus der kontinuierlichen Materieentstehung und der DSTT-Drift.

\subsection{Dichteverteilung}
\label{sub:kosmo_dichte}

Die Materieentstehungsrate im fraktalen Raum ist (siehe Anhang~\ref{app:materieentstehung}):

\[
\dot{\rho}_{\text{cre}}(r) \propto r^{D-3} = r^{-0,296}
\]

Im stationären Zustand führt dies über die Kontinuitätsgleichung zu einer radialen Dichteverteilung:

\[
\rho_{\text{Gal}}(r) \propto r^{-(3-D)} = r^{-0,296}
\]

\subsection{Rotationskurve}
\label{sub:kosmo_rotation}

Die kumulierte Masse innerhalb des Radius \( r \) ist:

\[
M(r) \propto \int_0^r \rho(r') r'^2 dr' \propto r^{2 + (3-D)} = r^{2,296}
\]

Daraus folgt die Rotationskurve:

\[
v(r) = \sqrt{\frac{GM(r)}{r}} \propto r^{0,648}
\]

Für \( D \approx 2,704 \) ergibt sich \( v(r) \propto r^{0,648} \) – dies ist eine fast flache Rotationskurve, konsistent mit Beobachtungen (\cite{Rubin1970}).

\section{Kontinuierliche Materieentstehung}
\label{sec:kosmo_materieentstehung}

Die Materieentstehung wird durch die negative Vakuumenergiedichte des fraktalen Raumes vermittelt. Aus Anhang~\ref{app:materieentstehung} folgt:

\[
\dot{\rho}_{\text{cre}}(r) = \kappa \cdot \frac{\hbar c}{l_0^4 c^2} \left( \frac{l_0}{r} \right)^{3-D} \propto r^{D-3}
\]

Die Materieentstehung ist also im Zentrum am stärksten und fällt mit einer gebrochenen Potenz ab, die direkt aus der fraktalen Dimension \( D \) folgt. Die vollständige Energiebilanz inklusive der negativen Vakuumenergiedichte findet sich in Anhang~\ref{app:materieentstehung}.

\subsection{Drift und Einfang}
\label{sub:kosmo_drift_einfang}

Die DSTT-Drift treibt die neu entstehende Materie nach außen. Die radiale Driftgeschwindigkeit wird in Anhang~\ref{app:identitaetsbeziehung} hergeleitet:

\[
v_{\text{drift}}(r) = r\dot{\phi} \sqrt{1-\beta}
\]

Für Kreisbahnen als Näherung (\(\beta \approx 1\)) folgt die Skalierung \( v_{\text{drift}} \propto r^{-1/2} \).

Die Einfangwahrscheinlichkeit durch den Zentralkörper ist (Herleitung in Anhang~\ref{app:identitaetsbeziehung}):

\[
\alpha_{\text{in}} = \left( \frac{r_c}{r_{\text{max}}} \right)^{3-D}
\]

Hierbei ist \( r_c \) der Einfangradius, innerhalb dessen neu entstehende Materie vom Zentralkörper absorbiert wird, und \( r_{\text{max}} \) der äußere Radius der Entstehungszone des Quantenpotentials. Für eine Galaxie mit zentralem BEC-Kern gilt \( r_c \sim 200 \, \text{m} \) (Oberfläche des BEC-Objekts, siehe Kapitel~\ref{ch:bec_objekte}) und \( r_{\text{max}} \sim 5 \times 10^{11} \, \text{m} \) (äußerer Rand der Entstehungszone), woraus sich \( \alpha_{\text{in}} \approx 0,001 \) ergibt.

\section{Galaxienwachstum}
\label{sec:kosmo_galaxienwachstum}

Während der Raum selbst nicht expandiert (das Universum ist stationär), wachsen gebundene Strukturen wie Galaxien durch kontinuierliche Materieentstehung und die DSTT-Drift. Aus der Drift-Gleichung \( dr/dt = v_{\text{drift}}(r) \) mit \( v_{\text{drift}} \propto r^{-1/2} \) folgt durch Integration:

\[
r(t) \propto t^{2/3}
\]

Die Materieentstehungsrate skaliert mit \( r^{3-D} \propto t^{2(3-D)/3} = t^{0,197} \). Daraus folgt für die kumulierte Masse der Galaxie:

\[
M_{\text{Gal}}(t) \propto t^{1,197}
\]

Der Exponent liegt nahe an 1, was ein nahezu lineares Wachstum der Galaxienmasse mit der Zeit bedeutet.

\subsection{Massenverhältnis Galaxie zu Kern}
\label{sub:kosmo_massenverhaeltnis}

Das Massenverhältnis von Galaxie zu zentralem BEC-Kern (siehe Kapitel~\ref{ch:bec_objekte}) ist:

\[
\frac{M_{\text{Gal}}}{M_{\text{Kern}}} = \frac{1 - \alpha_{\text{in}}}{\alpha_{\text{in}}} \approx 1000
\]

Mit \( M_{\text{Kern}} = M_{\text{BEC}} \approx 3 \times 10^6 M_\odot \) ergibt sich \( M_{\text{Gal}} \approx 3 \times 10^9 M_\odot \) – typisch für große Galaxien.

\subsection{Die Sonne als Mikrokosmos}
\label{sub:kosmo_sonne}

Das Modell der BEC-Objekte und der DSTT-Drift ist skalierbar. 
Überträgt man es auf die Masse der Sonne (\( M = M_{\odot} \)), 
so ergeben sich folgende charakteristische Skalen für das solare System:

\begin{align}
r_{\text{max}} &\sim 1,5 \times 10^{13} \, \text{m} \tag{14.20} \\
r_c &\sim 1,4 \times 10^{-11} \, \text{m} \tag{14.21} \\
\alpha_{\text{in}} &\approx 1,2 \times 10^{-7} \tag{14.22}
\end{align}

Dabei ist \( r_{\text{max}} \) der äußere Rand der Entstehungszone im solaren System, 
\( r_c \) der Einfangradius (in der Größenordnung der 
Compton-Wellenlänge des Protons) und \( \alpha_{\text{in}} \) die 
Einfangwahrscheinlichkeit neu entstehender Materie durch die Sonne.

Diese Werte sind spezifisch für die Skalierung auf die Sonnenmasse und folgen aus der fraktalen Geometrie des Raumes mit \( D \approx 2,704 \).

\subsubsection{Interpretation des Sonnenwinds}
\label{sub:kosmo_sonnenwind}

Im Standardmodell wird der Sonnenwind als Materieausstoß aus der 
Sonnenkorona interpretiert. Die WDBT+ bietet eine alternative 
Erklärung:

\begin{enumerate}
    \item Neue Materie entsteht im Quantenpotential des solaren 
          BEC-Objekts in der Nähe des Zentrums (\( r \sim 10^{-11} \, \text{m} \)).
    \item Die DSTT-Drift treibt diese Materie radial nach außen.
    \item Nur ein winziger Bruchteil (\( \alpha_{\text{in}} \approx 10^{-7} \)) 
          wird von der Sonne eingefangen.
    \item Der Großteil der Materie verlässt die Sonne als Sonnenwind.
\end{enumerate}

Eine wichtige Konsequenz: Die Sonne verliert durch diesen Prozess 
keine signifikante Masse, da der Sonnenwind nicht aus der Sonne 
selbst stammt, sondern aus neu entstehender Materie (siehe auch Kapitel~\ref{ch:testbare_vorhersagen}). 
Präzise Messungen der Sonnenmasse über längere Zeiträume 
sollten daher keinen Masseverlust zeigen – eine testbare 
Vorhersage der Theorie.

\subsubsection{Skalierbarkeit des Modells}
\label{sub:kosmo_skalierbarkeit}

Die fraktale Raumdimension \( D \approx 2,704 \) bestimmt die 
Exponenten aller Skalierungsgesetze. Für die Sonne (\( M = M_{\odot} \)) 
liefert das Modell konsistente Größenordnungen, die mit 
Beobachtungen des Sonnenwinds verträglich sind. Eine detaillierte 
quantitative Vorhersage erfordert die Kenntnis der genauen 
Dichteverteilung im solaren BEC-Objekt – eine Aufgabe für 
zukünftige Forschung.

\section{Teilchen als Interferenzmuster}
\label{sec:kosmo_teilchen_interferenz}

In der WDBT+ sind Teilchen keine fundamentalen Entitäten mit einem festen Informationswert. Vielmehr sind sie stabile, nichtlineare Interferenzmuster der komplexen Informationswellen \( I_k^{(n)} \). Die Stabilität dieser Muster wird durch die spezifischen Phasenbeziehungen der Wellen und die Randbedingungen des fraktalen Netzwerks gewährleistet. Die Amplitude eines Musters an einem Knoten \( k \) kann dabei Werte wie 0,01, 0,25 oder 1,0 annehmen, die in früheren Modellen fälschlicherweise als intrinsische Teilcheneigenschaften interpretiert wurden. Die wahre Natur eines Teilchens ist seine räumlich-zeitliche Kohärenzstruktur im Informationsfeld.

\section{Beobachtbare Signaturen: Der EHT-Schatten}
\label{sec:kosmo_eht}

Ein BEC-Objekt mit \( M = M_{\text{BEC}} \) hat \( R = l_0 \approx 10^{-15} \, \text{m} \) – viel zu klein, um direkt abgebildet zu werden. Was das Event Horizon Telescope (EHT) zeigt (\cite{EHT2019}, \cite{EHT2022}), ist der Einflussbereich: die Region, in der Licht stark abgelenkt wird (siehe Kapitel~\ref{ch:testbare_vorhersagen}).

Der Photonenorbit liegt bei:

\[
r_{\text{ph}} \approx \frac{3GM}{c^2}
\]

Für \( M = M_{\text{BEC}} \):

\[
r_{\text{ph}} \approx 1,8 \times 10^{10} \, \text{m} \approx 0,12 \, \text{AU}
\]

Der Schatten ist also makroskopisch. Testbare Unterschiede zur ART:

\begin{itemize}
    \item Feinstruktur des Schattens
    \item Polarisationsmuster
    \item Zeitliche Variabilität (Flares)
\end{itemize}

\section{Vergleich mit dem Standardmodell}
\label{sec:kosmo_vergleich}

Die folgende Tabelle fasst die Unterschiede zwischen der Standardkosmologie und der WDBT+ zusammen:

\begin{table}[htbp]
\centering
\caption{Gegenüberstellung der Standardkosmologie (\(\Lambda\)CDM) und der WDBT+.}
\begin{tabular}{lp{4.5cm}p{4.5cm}}
\hline
\textbf{Phänomen} & \textbf{Standardmodell (\(\Lambda\)CDM)} & \textbf{WDBT+} \\
\hline
Zentrale Objekte & Senken (verschlingen Materie) & Quellen (erzeugen Materie) \\
Materiefluss & Einwärts (Akkretion) & Auswärts (Drift) \\
Galaxien & Kollaps + Dunkle Materie & Wachstum von innen nach außen \\
Universum & Expandierend, Urknall & Stationär, ewig \\
Schwarze Löcher & Singularität + Horizont & BEC-Objekt, kein Horizont \\
Maximale Masse (Kern) & Keine natürliche Grenze & \( M_{\text{BEC}} \approx 3 \times 10^6 M_{\odot} \) \\
Super-massereich & Einzelnes Schwarzes Loch & BEC-Kern + Umgebung \\
Kollisionen & Immer Verschmelzung & Wie normale Objekte \\
Sonnenwind & Ausstoß aus der Korona & Neu entstehende Materie \\
Dunkle Materie & Benötigt & Nicht benötigt \\
Dunkle Energie & Benötigt & Nicht benötigt \\
\hline
\end{tabular}
\label{tab:kosmo_vergleich}
\end{table}

\section{Zusammenfassung}
\label{sec:kosmo_zusammenfassung}

Die WDBT+ führt zu einem stationären, nicht-expandierenden Universum:

\begin{itemize}
    \item Die Rotverschiebung ist keine Folge der Expansion, sondern kumulative Gravitation. Die allgemeine Form \( z = \frac{1}{c^2} \int \frac{GM(r)}{r^2} dr \) ist die Grundlage; die fraktale Dichteverteilung ist ein Spezialfall.
    \item Die effektive Hubble-Konstante \( H_{\text{eff}}(d) \) ist skalenabhängig: \( H_{\text{eff}}(d) \propto d^{0,704} \). Die Hubble-Spannung ist eine Konsequenz dieser Skalenabhängigkeit, kein Rätsel.
    \item Dieselbe Korrelationslänge \( L_{Q,0} \approx 2{,}0 \times 10^{46} \, \text{m} \) bestimmt sowohl die Neutrinomasse \( m_\nu \approx 0,1 \, \text{eV}/c^2 \) als auch die Skalenabhängigkeit von \( H_{\text{eff}} \) (siehe Anhang~\ref{app:neutrino_oszillation}). Dies ist ein zentrales konsistentes Element der Theorie.
    \item Flache Rotationskurven werden ohne Dunkle Materie durch das \( \Omega \)-Feld und die anisotrope Trägheit der DSTT erklärt.
    \item Das Quantenpotential erzeugt kontinuierlich neue Materie, die durch DSTT-Drift nach außen getrieben wird.
    \item Galaxien wachsen von innen nach außen; das Massenverhältnis Galaxie/Kern folgt aus der fraktalen Raumdimension \( D \).
    \item Teilchen sind stabile Interferenzmuster der komplexen Informationswellen im fraktalen Netzwerk.
    \item Der Sonnenwind wird durch Materieentstehung in Sonnennähe erklärt – die Sonne verliert keine signifikante Masse.
    \item Der EHT-Schatten ist makroskopisch, aber die Feinstruktur unterscheidet sich von ART-Vorhersagen.
\end{itemize}

Die WDBT+ kommt ohne Urknall, Dunkle Materie und Dunkle Energie aus – und liefert ein konsistentes, singularitätenfreies Bild des Kosmos.